% Part: lambda-calculus
% Chapter: lambda-definability
% Section: partial-recursive-functions

\documentclass[../../../include/open-logic-section]{subfiles}

\begin{document}

\olfileid[fr]{lam}{rep}{par}
\olsection{Fonctions récursives partielles \texorpdfstring{$\lambda$}{λ}-définissables}

Les fonctions récursives partielles
s'obtiennent à partir des fonctions
initiales par composition, récurrence
primitive et minimisation non bornée.
Elles se distinguent des fonctions
récursives générales par le fait que
les fonctions auxquelles on applique
la recherche non bornée n'ont pas à
être régulières. Sans cette exigence,
une fonction obtenue par minimisation
peut ne pas être définie sur certaines
entrées.

On pourrait croire que la méthode
employée pour les fonctions récursives
générales, qui sont totales, démontre
aussi que toutes les fonctions
récursives partielles sont
$\lambda$-définissables. Par exemple,
si $F$ et~$G$ définissent respectivement
$f$ et~$g$, le terme
$\lambd[x][F(Gx)]$ semble définir
leur composition. Lorsque $g(x)$
est indéfinie, toutefois, $Gx$
n'a pas de forme normale. Souvent,
$F(Gx)$ n'en a pas non plus, comme
il le faudrait; mais cela n'est
pas garanti. Prenons
\[
  F\ident\lambd[x][\num{0}].
\]
Même si $Gx$ n'a pas de forme normale,
$F(Gx)\red\num{0}$ en a une.

On peut représenter toute fonction récursive partielle par une
recherche non bornée qui ne comporte que des tests totaux. Le
théorème de forme normale de Kleene permet de choisir ces tests.

\begin{thm}
\ollabel{thm:partial-recursive-lambda}
Toute fonction récursive partielle est $\lambda$-définissable.
\end{thm}

\begin{proof}
Soit $f$ une fonction récursive partielle à $k$ arguments.
D'après le \olref[cmp][rec][nft]{thm:kleene-nf}, on peut écrire
\[
  f(\vec x)\simeq u(\umin{s}{R(\vec x,s)}),
\]
où $R$ est récursive primitive et $u$ est récursive primitive.
Pour $k>1$, on compose la relation de Kleene avec un codage
récursif primitif des tuples et on fixe l'indice de $f$.
Pour $k=0$, on fixe aussi une entrée factice. Soit $C$ un
lambda-terme qui définit la fonction récursive primitive
$1-\Char{R}$ et soit $U$ un lambda-terme qui définit~$u$.
Ces représentants existent par les constructions pour les
fonctions récursives primitives. Sur des numéraux de Church,
$C$ renvoie $\num{0}$ si le test réussit et $\num{1}$ sinon.

Utilisons le sélecteur $D(M,N,Z)=Z(K(N))M$ et le successeur
de l'introduction au lambda-calcul, ainsi que le combinateur
de Turing $Y$. Posons
\[
  L\ident\lambd[g][\lambd[\vec x\,s][
    D(U(s),g(\vec x,\fn{Succ}(s)),C(\vec x,s))]] .
\]
Le terme $\lambd[\vec x][(Y L)(\vec x,\num{0})]$
définit~$f$. Pour $k=0$, on omet l'abstraction et le passage
du tuple vide. L'écriture des termes est comprise à renommage
près des variables liées.

Vérifions la recherche. Si $Z\red\num{0}$, alors
$D(a,b,Z)\red a$ pour des variables fraîches $a,b$.
Si $Z\red\num{1}$, alors $D(a,b,Z)\red b$.
Le théorème de standardisation assure que la réduction externe
gauche atteint cette variable
(\href{https://lsv.ens-paris-saclay.fr/~goubault/Lambda/lambda.pdf\#page=20}%
{Goubault-Larrecq, \emph{Lambda-calcul et langages fonctionnels},
\S\,3.4, théorème~4}). Jusqu'à cette variable, la réduction
n'effectue que des pas de tête : une variable de tête et ses
arguments resteraient en place dans toute réduction ultérieure.
En remplaçant $a,b$ par des termes quelconques, ces mêmes pas
de tête conduisent à la branche choisie avant toute contraction
dans les branches.

Après le dépliage de $Y L$, ce calcul de tête choisit donc
$U(s)$ si le test réussit et l'appel suivant sinon. Au départ,
$s=\num{0}$. À chaque appel suivant, $s$ est remplacé par
$\fn{Succ}(s)$, qui se réduit au numéral suivant. Les
représentations de $C$ et de $U$ restent correctes sur ces
arguments, puisque la réduction se propage dans les applications.
Si un premier témoin $s_0$ existe, le terme atteint la branche
$U(s)$ avec $s\red\num{s_0}$, puis le numéral de $u(s_0)$.
Si aucun témoin n'existe, les dépliages et les sélections donnent
une réduction de tête infinie. Une telle réduction est le
préfixe de la stratégie externe gauche ; par standardisation,
le terme ne possède alors aucune forme normale.
\end{proof}

\end{document}
